% Glosarium. Setiap definisi diturunkan dari kalimat yang mendefinisikan
% istilah itu di babnya, dan istilahnya mengikuti daftar istilah di CLAUDE.md.
% Istilah baru yang dipakai naskah ditambahkan di sini pada perubahan yang sama.
\chapter{Glosarium}

\begin{description}
  \setlength{\itemsep}{2pt}
  \item[\textit{Activation cost}] Selisih \textit{startup time} antara
    \textit{eager activation} dan \textit{lazy activation} pada arsitektur
    \textit{plugin}.
  \item[\textit{Adapter}] Kelas yang memenuhi sebuah \textit{port} dengan
    teknologi tertentu, misalnya SQLite, sehingga \textit{core logic} tidak perlu
    mengenal teknologi itu.
  \item[\textit{Adapter swap cost}] Banyaknya \textit{file} dan baris yang
    berubah ketika teknologi di balik sebuah \textit{port} diganti.
  \item[\textit{Architecture sinkhole anti-pattern}] Keadaan ketika
    \textit{request} mengalir melewati beberapa \textit{layer} tanpa satu pun
    \textit{layer} mengerjakan sesuatu.
  \item[Arsitektur perangkat lunak] Susunan bagian besar sebuah sistem beserta
    aturan hubungan di antaranya. Pembedanya dari rancangan rinci adalah biaya
    untuk mengubahnya.
  \item[\textit{Backlog}] Tumpukan \textit{message} di antrean yang belum
    dikerjakan \textit{event processor}. Pada susunan berbasis event, waktu
    tunggu berpindah ke tumpukan ini.
  \item[\textit{Batch}] Cara menjalankan \textit{pipeline} yang menuntaskan satu
    \textit{stage} untuk seluruh \textit{file} sebelum \textit{stage} berikutnya
    mulai.
  \item[\textit{Binding}] Perantara berupa properti yang dapat diamati, yang
    menyalin perubahan antara ViewModel dan View. Nama lengkapnya \textit{data
    binding}.
  \item[\textit{Blocking}] Waktu yang dihabiskan \textit{publisher} untuk
    menunggu pihak yang dipanggilnya selesai sebelum dapat kembali.
  \item[\textit{Call stack}] Deretan pemanggilan metode yang sedang berlangsung,
    misalnya dari titik masuk \textit{request} sampai baris yang membaca basis
    data.
  \item[\textit{Centralized index}] Susunan \textit{peer-to-peer} yang menitipkan
    seluruh \textit{key} kepada satu \textit{node}, sehingga setiap
    \textit{lookup} cukup bertanya kepadanya.
  \item[\textit{Change cost}] Banyaknya \textit{file} dan baris yang berubah
    ketika satu \textit{subscriber} baru disisipkan.
  \item[Client-server] Gaya yang membagi sistem menjadi klien yang meminta dan
    server yang melayani. Data bersama dan aturan bisnis disimpan di server.
  \item[\textit{Closed layer}] \textit{Layer} yang wajib dilewati setiap
    \textit{request}, sehingga panggilan tidak boleh melompatinya.
  \item[\textit{Cold start}] Pemanggilan fungsi \textit{serverless} ketika belum
    ada contoh yang hidup, sehingga modulnya dimuat dan konstruktornya dijalankan
    lebih dahulu.
  \item[\textit{Composition root}] Satu-satunya tempat yang mengenal seluruh
    \textit{adapter}, lalu menyerahkannya ke \textit{core logic} sebagai
    argumen.
  \item[\textit{Concurrency limit}] Batas banyaknya contoh fungsi yang boleh
    berjalan bersamaan. \textit{Request} di atas batas itu harus menunggu.
  \item[\textit{Contract}] Bentuk yang ditulis \textit{core} dan wajib dipenuhi
    setiap \textit{plugin}.
  \item[\textit{Core}] Bagian kecil pada microkernel yang menulis
    \textit{contract} dan membaca \textit{manifest}, tanpa pernah menyebut nama
    satu \textit{plugin} pun.
  \item[\textit{Core independence}] Sifat \textit{core} atau \textit{core logic}
    yang tidak menyebut satu pun nama teknologi luar maupun nama penyumbangnya.
    Sifat itu diperiksa dengan menghitung penyebutannya.
  \item[\textit{Core logic}] Kode yang memuat aturan bisnis pada \textit{ports
    and adapters}, terpisah dari teknologi \textit{storage} dan antarmuka.
  \item[CQRS] Command Query Responsibility Segregation, yaitu pemisahan model
    yang dipakai menulis dari model yang dipakai membaca, dengan proyeksi
    sebagai penjembatannya.
  \item[\textit{Data collision}] Kejadian ketika dua unit mengubah baris yang
    sama di dalam \textit{replication lag}, sehingga satu perubahan menimpa
    perubahan lain.
  \item[\textit{Data grid}] Salinan data di dalam memori sebuah
    \textit{processing unit}, tempat unit itu menjawab \textit{request}.
  \item[\textit{Data loss}] Hilangnya perubahan yang sudah dijawab kepada
    pengguna tetapi belum tersimpan ke basis data ketika sebuah unit mati.
  \item[\textit{Data pump}] Saluran yang membawa perubahan dari
    \textit{processing unit} ke \textit{data writer}, sehingga basis data berada
    di luar jalur \textit{request}.
  \item[\textit{Data writer}] Komponen yang menuliskan perubahan dari
    \textit{data pump} ke basis data sekelompok demi sekelompok.
  \item[\textit{Decoupling}] Pelepasan ikatan antara pihak yang mengumumkan
    sebuah peristiwa dan pihak yang mengerjakan lanjutannya.
  \item[\textit{Dependency}] Hubungan ketika sebuah kelas atau modul mengenal
    kelas atau modul lain, sehingga perubahan pada yang dikenal dapat memaksa
    yang mengenal ikut berubah.
  \item[\textit{Dependency inversion}] Pembalikan arah \textit{dependency},
    sehingga panahnya berangkat dari \textit{adapter} menuju \textit{core
    logic}.
  \item[\textit{Eager activation}] Cara mengaktifkan \textit{plugin} yang
    mengimpor seluruh modul saat aplikasi mulai.
  \item[\textit{Error isolation}] Kemampuan \textit{core} tetap berjalan ketika
    satu \textit{plugin} gagal, karena pemanggilannya dilindungi \textit{error
    handler}.
  \item[\textit{Event log}] Deretan event yang disimpan pada \textit{event
    sourcing}. Isinya mencatat yang terjadi, bukan yang diminta.
  \item[\textit{Event processor}] Komponen yang menerima event lalu mengerjakan
    satu langkah lanjutan, misalnya \textit{recorder}, \textit{validator}, atau
    \textit{notifier}.
  \item[\textit{Event sourcing}] Cara menyimpan yang hanya mencatat event, lalu
    menghitung ulang keadaan sekarang dari event tersebut.
  \item[\textit{Event-driven architecture}] Gaya yang membuat \textit{publisher}
    mengumumkan bahwa sesuatu telah terjadi lalu berhenti, sedangkan
    \textit{subscriber} yang berkepentingan mengerjakan lanjutannya.
  \item[\textit{Eventual consistency}] Keadaan ketika model baca menyusul model
    tulis sesudah selang tertentu, sehingga pembacaan sesaat sesudah penulisan
    dapat memuat angka lama.
  \item[\textit{Extension point}] Sebutan Eclipse untuk \textit{contract}, yaitu
    slot yang disediakan sebuah \textit{plugin} agar \textit{plugin} lain dapat
    memenuhinya.
  \item[\textit{Failure containment}] Banyaknya \textit{request} yang tetap
    tuntas ketika satu \textit{event processor} gagal.
  \item[\textit{Failure resilience}] Banyaknya \textit{lookup} yang masih
    terjawab ketika satu \textit{node} keluar dari jaringan.
  \item[\textit{Falsifiable claim}] Klaim yang dapat dibuktikan salah oleh satu
    pengamatan, misalnya klaim nol impor yang gugur oleh satu baris impor.
  \item[\textit{Fan-out}] Penggandaan satu \textit{stream} menjadi beberapa
    \textit{stream} yang berjalan sendiri-sendiri. Lawannya adalah
    \textit{merge}.
  \item[\textit{Fat controller}] \textit{Symptom} ketika Controller menyimpan
    aturan bisnis. Bentuk sejenisnya adalah \textit{fat model} dan \textit{fat
    view}.
  \item[\textit{Filter}] \textit{Stage} pada \textit{pipe-and-filter} yang
    menerima satu \textit{stream} lalu menghasilkan \textit{stream} baru, tanpa
    mengenal pengirim maupun penerimanya.
  \item[\textit{Filter independence}] Sifat \textit{filter} yang tidak menyebut
    nama \textit{filter} lain, sehingga urutannya tidak terkunci.
  \item[\textit{Finger table}] Daftar jalan pintas pada setiap \textit{node}
    Chord, sehingga \textit{lookup} dapat melompat jauh di \textit{identifier
    ring}.
  \item[\textit{First output latency}] Waktu sejak \textit{pipeline} mulai
    sampai \textit{file} keluaran pertama selesai.
  \item[Gerbang] Layanan pada mikroservis yang menerima \textit{request} dari
    luar, memanggil layanan lain, lalu menggabungkan jawabannya dengan kode
    penggabung.
  \item[\textit{Hexagonal architecture}] Nama lain \textit{ports and
    adapters}.
  \item[\textit{Hop}] Satu perpindahan panggilan, misalnya antar \textit{layer}
    di dalam satu proses atau antar layanan lewat jaringan.
  \item[\textit{Identifier ring}] Lingkaran bilangan tempat \textit{node} dan
    \textit{key} pada Chord diberi \textit{identifier}, sehingga letak sebuah
    \textit{key} dapat dihitung.
  \item[\textit{Idle cost}] Biaya yang tetap dibayar ketika \textit{request}
    tidak datang, misalnya memori dan modul yang dipegang layanan yang hidup
    terus.
  \item[\textit{Inbound port}] \textit{Port} yang dipakai pihak luar untuk
    memanggil \textit{core logic}, misalnya oleh antarmuka pengguna.
  \item[\textit{Latency cost}] Tambahan latensi ujung ke ujung yang dibayar
    sebuah susunan dibanding susunan sinkron yang mengerjakan pekerjaan sama.
  \item[\textit{Layer}] Satu tingkat pada \textit{layered architecture} yang
    mengerjakan satu jenis urusan teknis, misalnya presentasi atau persistensi.
  \item[\textit{Layer violation}] Rujukan dari sebuah kelas ke kelas yang tidak
    berada tepat satu \textit{layer} di bawahnya.
  \item[\textit{Layered architecture}] Gaya yang menata aplikasi menjadi
    \textit{layer} bertumpuk, dengan panggilan yang hanya boleh mengarah ke
    bawah.
  \item[\textit{Lazy activation}] Cara mengaktifkan \textit{plugin} yang membaca
    \textit{manifest} saat mulai, lalu menunda impor modul sampai perintahnya
    dipanggil.
  \item[\textit{Lookup}] Pekerjaan menemukan \textit{node} yang menyimpan sebuah
    \textit{key}.
  \item[\textit{Lookup cost}] Banyaknya \textit{message} atau \textit{hop} yang
    dihabiskan satu \textit{lookup}.
  \item[\textit{Manifest}] \textit{File} terpisah dari kode yang menyatakan
    kemampuan sebuah \textit{plugin}, sehingga daftar perintah dapat dibaca
    tanpa mengimpor modulnya.
  \item[Median] Nilai tengah hasil pengukuran sesudah diurutkan. Separuh hasil
    berada di bawahnya dan separuh di atasnya.
  \item[\textit{Merge}] Penyatuan beberapa \textit{stream} menjadi satu
    \textit{stream}.
  \item[Mesin orkestrasi] Komponen yang membaca urutan langkah alur kerja,
    memanggil layanan satu per satu, lalu menyimpan keadaan proses sesudah tiap
    langkah.
  \item[\textit{Message scalability}] Pertumbuhan jumlah \textit{message} ketika
    jumlah \textit{node} dinaikkan, dibanding pertumbuhan jumlah \textit{node}
    itu sendiri.
  \item[Microkernel] Gaya yang membagi sistem menjadi \textit{core} kecil
    beserta \textit{plugin} yang dapat dipasang dan dicabut.
  \item[Mikroservis] Gaya yang membagi sistem menjadi banyak layanan kecil,
    masing-masing dengan proses dan basis datanya sendiri.
  \item[Model-View-*] Kelompok pola yang memisahkan tampilan dari aturan bisnis.
    Variannya Model-View-Controller (MVC), Model-View-Presenter (MVP),
    Model-View-ViewModel (MVVM), dan Model-View-Intent (MVI).
  \item[Monolit] Aplikasi yang berjalan sebagai satu proses dengan satu basis
    data, dipakai sebagai pembanding pada Bab~11 dan Bab~13.
  \item[\textit{Neighbor}] \textit{Node} yang dikenal langsung oleh sebuah
    \textit{node}, sehingga pertanyaan dapat dikirim kepadanya tanpa perantara.
  \item[\textit{Node}] Satu peserta jaringan \textit{peer-to-peer} yang
    berkedudukan sama dengan peserta lain.
  \item[\textit{Observer}] Mekanisme pemberitahuan yang membuat Model tidak perlu
    mengetahui jenis View yang mengamatinya.
  \item[\textit{Open layer}] \textit{Layer} yang boleh dilompati \textit{request}
    tanpa aturan bisnis, sehingga jalannya berkurang satu \textit{hop}.
  \item[\textit{Orchestration-driven service-oriented architecture}] Gaya yang
    menata layanan menjadi beberapa \textit{layer}, lalu merangkainya dengan
    satu mesin orkestrasi.
  \item[\textit{Outbound port}] \textit{Port} yang dipakai \textit{core logic}
    untuk meminta sesuatu dari luar, misalnya membaca data.
  \item[\textit{Pass-through}] Metode yang tubuhnya hanya satu baris
    \texttt{return} berisi pemanggilan, tanpa menambahkan apa pun.
  \item[\textit{Passive view}] Bentuk MVP yang View-nya sama sekali tidak
    mengenal komponen lain.
  \item[\textit{Peer-to-peer}] Gaya yang menata sistem menjadi \textit{node}
    berkedudukan sama, sehingga tidak ada pihak yang selalu melayani.
  \item[Persentil ke-95] Nilai yang tidak dilampaui oleh 95 persen hasil
    pengukuran. Nilai ini memperlihatkan ekor sebaran yang tidak tampak pada
    median.
  \item[\textit{Pipe}] Penyambung antar \textit{stage} pada
    \textit{pipe-and-filter}.
  \item[\textit{Pipe-and-filter}] Gaya yang menata pekerjaan menjadi deretan
    \textit{stage} yang masing-masing mengerjakan satu perubahan kecil, disambung
    oleh \textit{pipe}.
  \item[\textit{Pipeline}] Rangkaian \textit{stage} yang dilewati data secara
    berurutan.
  \item[\textit{Plugin}] Bagian tambahan yang memenuhi \textit{contract} milik
    \textit{core}, ditulis pihak lain, dan dapat dipasang atau dicabut.
  \item[\textit{Port}] Antarmuka yang menyebut \textit{behavior} yang dibutuhkan
    \textit{core logic}, tanpa menyebut teknologi yang menyediakannya.
  \item[\textit{Ports and adapters}] Gaya yang membalik arah \textit{dependency},
    sehingga \textit{core logic} menyatakan kebutuhannya sebagai \textit{port}
    lalu \textit{adapter} memenuhinya.
  \item[\textit{Processing unit}] Unit pada \textit{space-based architecture}
    yang memuat modul aplikasi beserta \textit{data grid}-nya.
  \item[Proyeksi] Fungsi yang melipat satu perubahan ke atas ringkasan yang sudah
    ada. Pada CQRS, proyeksi menjembatani model tulis dan model baca.
  \item[\textit{Publisher}] Pihak yang mengumumkan bahwa sesuatu telah terjadi,
    lalu berhenti tanpa menyebut pihak yang mengerjakan lanjutannya.
  \item[\textit{Query flooding}] Cara \textit{lookup} pada susunan
    \textit{unstructured} yang meneruskan pertanyaan ke seluruh
    \textit{neighbor}, lalu diteruskan lagi oleh setiap \textit{node} yang
    menerimanya.
  \item[\textit{Range}] Selisih antara putaran tercepat dan putaran terlambat
    pada satu pengukuran berulang.
  \item[\textit{Read cost}] Biaya menjawab satu kueri.
  \item[\textit{Registry}] Simpanan \textit{manifest} yang sudah terbaca,
    sehingga \textit{core} dapat mencari \textit{plugin} tanpa memindai ulang.
  \item[\textit{Replay cost}] Biaya melipat ulang seluruh \textit{event log}
    untuk menghitung keadaan sekarang.
  \item[\textit{Replication lag}] Selang sejak satu unit mengubah data di
    \textit{data grid}-nya sampai perubahan itu terpakai unit lain.
  \item[\textit{Scale to zero}] Turunnya jumlah contoh fungsi yang hidup menjadi
    nol ketika \textit{request} tidak datang.
  \item[\textit{Service-based architecture}] Gaya yang memecah aplikasi menjadi
    sedikit layanan domain berbutir kasar, lazimnya empat sampai dua belas,
    dengan satu basis data bersama.
  \item[\textit{Serverless}] Gaya yang menyerahkan urusan menghidupkan dan
    menjaga proses kepada penyedia, dengan satu fungsi sebagai satuan
    penyebaran.
  \item[Snapshot] Ringkasan keadaan pada satu titik di \textit{event log},
    sehingga penghitungan ulang cukup dimulai dari titik itu.
  \item[\textit{Space-based architecture}] Gaya yang menyalin data yang
    dibutuhkan setiap \textit{request} ke memori setiap \textit{processing
    unit}, sehingga basis data keluar dari jalur \textit{request}.
  \item[\textit{Stage}] Satu langkah pada \textit{pipeline} yang mengerjakan
    satu perubahan kecil.
  \item[\textit{Stage reordering}] Penukaran urutan dua \textit{stage}.
    Penukaran itu sah bila bentuk datanya sama, tetapi belum tentu menghasilkan
    \textit{file} yang sama.
  \item[\textit{Stale read}] Jawaban dengan angka lama yang diberikan sebuah unit
    selama \textit{replication lag} belum lewat.
  \item[\textit{Startup time}] Waktu sejak aplikasi dijalankan sampai siap
    menerima perintah.
  \item[\textit{Stream}] Aliran data yang diterima dan dihasilkan sebuah
    \textit{filter}.
  \item[\textit{Streaming}] Cara menjalankan \textit{pipeline} yang membawa satu
    \textit{file} menempuh seluruh \textit{stage} sebelum \textit{file}
    berikutnya dibuka.
  \item[\textit{Structured}] Susunan \textit{peer-to-peer} yang menghitung letak
    sebuah \textit{key} dari \textit{key}-nya sendiri, misalnya Chord.
  \item[\textit{Stub}] Pengganti \textit{adapter} untuk pengujian yang
    mengembalikan nilai tetap tanpa membuka koneksi basis data.
  \item[\textit{Subscriber}] Pihak yang mendaftarkan dirinya pada sebuah
    \textit{topic}, lalu mengerjakan lanjutan setiap event pada \textit{topic}
    itu.
  \item[\textit{Successor}] \textit{Node} pertama yang \textit{identifier}-nya
    sama dengan atau sesudah \textit{identifier} sebuah \textit{key}. Pada
    Chord, \textit{node} itulah yang menyimpan \textit{key} tersebut.
  \item[\textit{System of record}] Basis data yang memegang catatan resmi,
    sehingga salinan di memori unit boleh dibuang lalu dibentuk ulang darinya.
  \item[\textit{Temporal query}] Pertanyaan tentang keadaan pada satu titik
    waktu di masa lampau, misalnya saldo pada tanggal dan jam tertentu.
  \item[Terukur] Sifat klaim arsitektur yang menyebut besaran beserta
    satuannya, berasal dari program yang dijalankan menurut prosedur tertulis,
    dan dapat diperoleh ulang di komputer lain.
  \item[\textit{Testing in isolation}] Pengujian \textit{core logic} tanpa basis
    data, dengan \textit{stub} sebagai pengganti \textit{adapter}.
  \item[\textit{Throughput scaling}] Pertambahan \textit{request} yang
    terlayani per detik ketika jumlah \textit{processing unit} ditambah.
  \item[\textit{Topic}] Nama saluran tempat \textit{publisher} mengumumkan event
    dan \textit{subscriber} mendaftarkan dirinya.
  \item[Topologi broker] Susunan berbasis event tanpa pengatur urutan. Setiap
    \textit{subscriber} menerbitkan event baru sesudah mengerjakan bagiannya.
  \item[Topologi mediator] Susunan berbasis event dengan satu mediator yang
    memegang daftar langkah, memanggil \textit{event processor} satu per satu,
    lalu menunggu balasannya.
  \item[\textit{Trade-off}] Pertukaran ketika satu sifat sistem diperoleh dengan
    mengorbankan sifat lain.
  \item[\textit{Tuple space}] Ruang bersama tempat unit bertukar data. Gagasan
    itu menjadi asal nama \textit{space-based architecture}.
  \item[\textit{Unstructured}] Susunan \textit{peer-to-peer} yang menyebarkan
    \textit{key} ke sembarang \textit{node}, sehingga tidak ada yang tahu
    letaknya.
  \item[\textit{Warm start}] Pemanggilan fungsi \textit{serverless} pada contoh
    yang sudah hidup, sehingga modulnya tidak dimuat ulang.
  \item[\textit{Win count}] Banyaknya putaran tempat satu pilihan lebih cepat
    daripada pilihan lain pada pengukuran berulang.
  \item[\textit{Write contention}] Perebutan giliran menulis ketika seluruh unit
    menulis ke basis data yang sama.
  \item[\textit{Write cost}] Biaya menjalankan satu perintah, misalnya latensi
    tulis pada CQRS.
  \item[\textit{Write-behind}] Cara menulis yang mengumpulkan perubahan, lalu
    mengirimnya ke \textit{storage} sebagai satu operasi kelompok.
\end{description}
